% !PRINTED 183
\weberpage{183}

%% [running head, not body text -- commented out 2026-08-07] Page furniture (running
%% head, folio, signature) is excluded from this edition by standing policy.
%% \noindent \S.\,47. \hfill \textit{Normalformen.} \hfill 183

%<<CONT
setzen, so ergiebt die Vergleichung von (6) mit (9):
\weq{10}{\alpha_{k,i} = \sum^s \frac{\varphi_i'(A_s)\,\Phi'(\Theta) - \Phi'(A_s)\,\varphi_i'(\Theta)}{\Phi'(\Theta)}\, B_s'(y_k).}

Die rechte Seite dieser Gleichungen muss sich also auf eine Constante reduciren, und wir können ihren Werth finden, wenn wir für die $y$ irgend ein specielles Werthsystem setzen, das nur an die eine Bedingung gebunden ist, dass $\Phi'(\Theta)$ nicht verschwindet. Für dieses specielle Werthsystem sind die Werthe der $A_i$ durch die Gleichungen (5) bestimmt, und $\Theta$ ist bekannt, wenn wir für dies specielle Werthsystem der $A_i$ das Formenproblem (3) als gelöst voraussetzen. Dann sind durch (10) die Co\"efficienten $\alpha_{k,i}$ und damit die Substitution $L$ vollständig bestimmt.

\noindent\hrulefill
